% GENERATED FILE. Edit canonical lessons, then run npm run generate:books.
% canonical-source-hash: fnv1a64:170daa02678af322
% canonical-lessons: UR-C222-mazedar, UR-C222-bekar, UR-C222-anokha, UR-C222-naqli, UR-C222-sachcha

\chapter{Tasty, Useless, Unusual, Fake, Truthful}
\label{ch:ur-tasty-useless-unusual-fake-truthful}
\hlchaptermodality{222}

\begin{chapteropening}
\textbf{By the end of this chapter:} \emph{I can say tasty, useless, unusual, fake and truthful.}

Complete the last lesson of chapter 222: I can say tasty, useless, unusual, fake and truthful.
\end{chapteropening}

\section[\texorpdfstring{\ur{مزیدار} (mazedār)}{مزیدار (mazedār)}]{\ur{مزیدار} (mazedār) — tasty}

\label{lesson:UR-C222-mazedar}



\begin{quote}
Before the new one: say the Urdu for quiet, then the Urdu for splendid.
\end{quote}

\subsection*{You'll want to know: \ur{مزیدار}}
\textbf{\ur{مزیدار}} — \emph{mazedār} — \textquotedblleft{}tasty\textquotedblright{}.

\begin{cousinweb}
\textbf{\ur{مزیدار}} (\emph{mazedār}) is \textquotedblleft{}holding taste\textquotedblright{}: \emph{maza}, \textquotedblleft{}taste, relish\textquotedblright{}, usually traced to Persian \emph{maza}, with the Persian \emph{-dār}, \textquotedblleft{}holder\textquotedblright{}.
\end{cousinweb}

\begin{grammarlens}[title={What you've built}]
One more word for this chapter.
\end{grammarlens}

\subsection*{Your turn}
\begin{itemize}
\raggedright
  \item \emph{Say it:} \emph{mazedār}
  \item \emph{Say it:} \emph{mazedār}, once more
  \item \emph{Say it:} say \emph{mazedār}
  \item \emph{Recall:} say the Urdu for quiet, then the Urdu for splendid, then say \emph{mazedār} again
\end{itemize}

\begin{tcolorbox}[breakable,colback=teal!4,colframe=teal!35!black,title={Before you move on}]
What does \ur{مزیدار} mean? (\textquotedblleft{}Tasty\textquotedblright{}.) Say it once more.
\end{tcolorbox}
\section[\texorpdfstring{\ur{بیکار} (bekār)}{بیکار (bekār)}]{\ur{بیکار} (bekār) — useless}

\label{lesson:UR-C222-bekar}



\begin{quote}
Before the new one: say the Urdu for splendid, then the Urdu for tasty.
\end{quote}

\subsection*{You'll want to know: \ur{بیکار}}
\textbf{\ur{بیکار}} — \emph{bekār} — \textquotedblleft{}useless\textquotedblright{}.

\begin{cousinweb}
\textbf{\ur{بیکار}} (\emph{bekār}) is a loan, usually traced to Persian \emph{bēkār}, \textquotedblleft{}idle, useless\textquotedblright{}. It is built from \emph{bē-}, \textquotedblleft{}without\textquotedblright{}, and \emph{kār}, \textquotedblleft{}work\textquotedblright{}.
\end{cousinweb}

\begin{grammarlens}[title={What you've built}]
One more word for this chapter.
\end{grammarlens}

\subsection*{Your turn}
\begin{itemize}
\raggedright
  \item \emph{Say it:} \emph{bekār}
  \item \emph{Say it:} \emph{bekār}, once more
  \item \emph{Say it:} say \emph{bekār}
  \item \emph{Recall:} say the Urdu for splendid, then the Urdu for tasty, then say \emph{bekār} again
\end{itemize}

\begin{tcolorbox}[breakable,colback=teal!4,colframe=teal!35!black,title={Before you move on}]
What does \ur{بیکار} mean? (\textquotedblleft{}Useless\textquotedblright{}.) Say it once more.
\end{tcolorbox}
\section[\texorpdfstring{\ur{انوکھا} (anokhā)}{انوکھا (anokhā)}]{\ur{انوکھا} (anokhā) — unusual}

\label{lesson:UR-C222-anokha}



\begin{quote}
Before the new one: say the Urdu for tasty, then the Urdu for useless.
\end{quote}

\subsection*{You'll want to know: \ur{انوکھا}}
\textbf{\ur{انوکھا}} — \emph{anokhā} — \textquotedblleft{}unusual\textquotedblright{}.

\begin{cousinweb}
\textbf{\ur{انوکھا}} (\emph{anokhā}) is a tadbhava, a word inherited from Sanskrit and reshaped by sound change. It comes from \emph{anapekṣa}, \textquotedblleft{}careless\textquotedblright{} (CDIAL 289). It is built on \emph{apekṣā}, \textquotedblleft{}consideration, expectation\textquotedblright{}: what is strange is what nobody expected.
\end{cousinweb}

\begin{grammarlens}[title={What you've built}]
One more word for this chapter.
\end{grammarlens}

\subsection*{Your turn}
\begin{itemize}
\raggedright
  \item \emph{Say it:} \emph{anokhā}
  \item \emph{Say it:} \emph{anokhā}, once more
  \item \emph{Say it:} say \emph{anokhā}
  \item \emph{Recall:} say the Urdu for tasty, then the Urdu for useless, then say \emph{anokhā} again
\end{itemize}

\begin{tcolorbox}[breakable,colback=teal!4,colframe=teal!35!black,title={Before you move on}]
What does \ur{انوکھا} mean? (\textquotedblleft{}Unusual\textquotedblright{}.) Say it once more.
\end{tcolorbox}
\section[\texorpdfstring{\ur{نقلی} (naqlī)}{نقلی (naqlī)}]{\ur{نقلی} (naqlī) — fake}

\label{lesson:UR-C222-naqli}



\begin{quote}
Before the new one: say the Urdu for useless, then the Urdu for unusual.
\end{quote}

\subsection*{You'll want to know: \ur{نقلی}}
\textbf{\ur{نقلی}} — \emph{naqlī} — \textquotedblleft{}fake\textquotedblright{}.

\begin{cousinweb}
\textbf{\ur{نقلی}} (\emph{naqlī}) is \emph{naql}, \textquotedblleft{}a copy\textquotedblright{}, with the ending \emph{-ī}, \textquotedblleft{}of\textquotedblright{}. \emph{naql} is Arabic, \textquotedblleft{}a moving across; a copying\textquotedblright{}, and reached Urdu through Persian (Lane, root \emph{n-q-l}).
\end{cousinweb}

\begin{grammarlens}[title={What you've built}]
One more word for this chapter.
\end{grammarlens}

\subsection*{Your turn}
\begin{itemize}
\raggedright
  \item \emph{Say it:} \emph{naqlī}
  \item \emph{Say it:} \emph{naqlī}, once more
  \item \emph{Say it:} say \emph{naqlī}
  \item \emph{Recall:} say the Urdu for useless, then the Urdu for unusual, then say \emph{naqlī} again
\end{itemize}

\begin{tcolorbox}[breakable,colback=teal!4,colframe=teal!35!black,title={Before you move on}]
What does \ur{نقلی} mean? (\textquotedblleft{}Fake\textquotedblright{}.) Say it once more.
\end{tcolorbox}
\section[\texorpdfstring{\ur{سچا} (sachchā)}{سچا (sachchā)}]{\ur{سچا} (sachchā) — truthful}

\label{lesson:UR-C222-sachcha}



\begin{quote}
Before the new one: say the Urdu for unusual, then the Urdu for fake.
\end{quote}

\subsection*{You'll want to know: \ur{سچا}}
\textbf{\ur{سچا}} — \emph{sachchā} — \textquotedblleft{}truthful\textquotedblright{}.

\begin{cousinweb}
\textbf{\ur{سچا}} (\emph{sachchā}) is a tadbhava, a word inherited from Sanskrit and reshaped by sound change. It comes from \emph{satya}, \textquotedblleft{}true\textquotedblright{} (CDIAL 13112).
\end{cousinweb}

\begin{grammarlens}[title={What you've built}]
One more word for this chapter.
\end{grammarlens}

\subsection*{Your turn}
\begin{itemize}
\raggedright
  \item \emph{Say it:} \emph{sachchā}
  \item \emph{Say it:} \emph{sachchā}, once more
  \item \emph{Say it:} say \emph{sachchā}
  \item \emph{Recall:} say the Urdu for unusual, then the Urdu for fake, then say \emph{sachchā} again
\end{itemize}

\begin{tcolorbox}[breakable,colback=teal!4,colframe=teal!35!black,title={Before you move on}]
What does \ur{سچا} mean? (\textquotedblleft{}Truthful\textquotedblright{}.) Say it once more.
\end{tcolorbox}
